% Lösungen Einheit 3 von 3 – Produktregel (zusammenbau v0.1)
\einheitenkopf{Einheit 3 von 3 · Produktregel}

\erg{17}{a) nur Produktregel – $e^x$ ist nicht verkettet. \quad b) $f'(x) = 1 \cdot e^x + x \cdot e^x = (x + 1) \cdot e^x$ \quad c) $f'(t) = 5 \cdot e^{-\frac{t}{20}} + 5t \cdot (-\frac{1}{20}) \cdot e^{-\frac{t}{20}} = 5 \cdot (1 - \frac{t}{20}) \cdot e^{-\frac{t}{20}}$ \quad d) $f'(x) = 3 \cdot 2x \cdot e^x + 3 \cdot (x^2 - 4) \cdot e^x = 3 \cdot (x^2 + 2x - 4) \cdot e^x$ \quad e) $f'(x) = (2x - 4) \cdot e^{-x} - (x^2 - 4x + 1) \cdot e^{-x} = (-x^2 + 6x - 5) \cdot e^{-x} = -(x - 1) \cdot (x - 5) \cdot e^{-x}$; mögliche Extremstellen: $x_1 = 1$, $x_2 = 5$ (Nullstellen des Polynomfaktors, $e^{-x} \neq 0$) \quad f) $f''(x) = -2x \cdot e^{-x} + (-x^2 - 1) \cdot (-1) \cdot e^{-x} = (x^2 - 2x + 1) \cdot e^{-x} = (x - 1)^2 \cdot e^{-x}$; wegen $(x - 1)^2 \geq 0$ und $e^{-x} > 0$ ist $f''(x) \geq 0$ – die Nullstelle $x = 1$ ist doppelt, ohne Vorzeichenwechsel: keine Wendestelle, der Graph ist überall linksgekrümmt. \quad g) $h'(2) = f'(2) \cdot g(2) + f(2) \cdot g'(2) = 0 \cdot 3 + (-1) \cdot (-1) = 1$ \quad h) $g'(x) = f'(x) \cdot e^{-x} - f(x) \cdot e^{-x} = (f'(x) - f(x)) \cdot e^{-x}$; waagerechte Tangente: $g'(a) = 0$; weil $e^{-a} \neq 0$, folgt $f'(a) - f(a) = 0$, also $f'(a) = f(a)$.}
\erg{18}{a) Ole hat nur die Faktoren einzeln abgeleitet und multipliziert; richtig nach der Produktregel: $f'(x) = 4x^3 \cdot e^x + x^4 \cdot e^x = (x^4 + 4x^3) \cdot e^x$. \quad b) Ein Produkt ist nur null, wenn ein Faktor null ist; $e^x$ ist nie null, also bleibt nur $x - 7 = 0$.}
